\documentclass[12pt,a4paper]{article}
\usepackage[utf8]{inputenc}
\usepackage{xeCJK}
\usepackage{geometry}
\usepackage{enumitem}
\usepackage{booktabs}
\usepackage{amsmath}
\usepackage{amssymb}
\usepackage{titlesec}
\usepackage{array}
\usepackage{multirow}
\usepackage{tikz}

\setmainfont{Times New Roman}
\setCJKmainfont{Microsoft JhengHei}

\geometry{margin=2.3cm, top=2.3cm, bottom=2.3cm}

% 行距設為 1.7 倍，使內容較為稀疏，較易達到 5 頁
\renewcommand{\baselinestretch}{1.7}

\titleformat{\section}
  {\normalfont\Large\bfseries}
  {\thesection}{1em}{}

\titleformat{\subsection}
  {\normalfont\large\bfseries}
  {\thesubsection}{1em}{}

\setlength{\parindent}{0pt}
\setlength{\parskip}{0.8em}

\begin{document}

\maketitle

\noindent\textbf{來源：} \texttt{300字專業摘要.docx} \quad
\textbf{格式範本：} \texttt{A2Z outline example.pdf}

\vspace{1em}

\section{Abstract}
本研究針對臺灣「開放山林」政策後高山登山事故倍增的議題，提出 TrailGuard TW，一個以證據為基礎的高山風險決策支援代理人。依 A2Z Outline 結構，本文從動機（A）、挑戰（B）、解藥（C）、方法設計（D）、實驗（E）與結果（F）六個面向，呈現研究概要。摘要內容摘錄自 \texttt{300字專業摘要.docx}，格式參照 \texttt{A2Z outline example.pdf}。

\subsection*{A. Attention Getter / Motivation（動機及引題）}
2019 年「開放山林」政策實施後，臺灣高山型國家公園入園人次減半，然而事故率卻倍增至 0.08–0.09\%。這顯示出山林資源開放與登山安全管理之間存在顯著的落差。隨著越來越多民眾進入高山地區，行前決策的品質成為影響登山安全的關鍵因素。

\subsection*{B. But（挑戰）}
全國山域事故中，迷路事故占 38.61\%，自組隊伍占 55.85\%。這兩項數據指出，許多登山事故並非源於不可抗力的天候或地形，而是源於行前缺乏完善的風險評估與決策支持。現有的登山資訊分散於不同平台，時效性不足，登山者難以快速取得整合性的風險資訊，導致行前決策品質不足。

\subsection*{C. Cure（解藥）}
本研究提出 TrailGuard TW，一個以證據為基礎（evidence-based）的高山風險決策支援代理人。系統將大語言模型（LLM）的語意理解能力與 Python 引擎的定量計算能力相結合，並透過時序知識圖譜管理資訊效期，以 Evidence Ledger 落實「無來源，無數字」的原則，並設置硬性紅旗機制直接輸出 NO-GO，確保高風險方案不會被誤判為安全。

\subsection*{D. Development（我們的方法設計）}
TrailGuard TW 的架構包含四個主要組件：（1）LLM 負責理解使用者需求並生成自然語言解釋；（2）Python 引擎計算 P50／P90 步程、天氣窗格與高山症風險；（3）時序知識圖譜管理資訊的效期與時間屬性；（4）Evidence Ledger 記錄每項建議的來源、時間戳與有效期限，硬性紅旗機制則在檢測到關鍵風險條件時直接輸出 NO-GO。

\subsection*{E. Experiments（實驗）}
研究以 120 個登山情境案例進行比較對照與消融組實驗。案例涵蓋不同天候、路線難度、海拔與行程長度，用於驗證系統在多約束規劃任務下的決策可靠性與風險識別表現。

\subsection*{F. Findings（結果）}
實驗以錯誤安全率（False-Safe Rate, FSR）作為核心評估指標，預期 TrailGuard TW 能顯著降低高風險登山方案被誤判為安全的比例，並建立可稽核、由人類做最終決策的高風險領域代理人範式。

\section{Introduction}
本章節依 A2Z Outline 的 G–L 結構，詳述研究背景、文獻現況與研究貢獻。

\subsection*{G. General 現況}
臺灣擁有豐富的高山資源與多樣的登山環境。2019 年「開放山林」政策推動後，公眾進入高山型國家公園的門檻降低，登山活動愈發普及。然而，隨著入山人口增加，高山登山事故亦隨之攀升，成為戶外安全管理面臨的重要課題。

\subsection*{H. However 挑戰}
根據本研究統計，全國山域事故中迷路占比達 38.61\%，自組隊伍占比更高達 55.85\%。這兩項數據顯示，許多登山事故並非無法預見的意外，而是與行前決策品質不足、風險意識薄弱相關。現有登山資訊分散於不同網站與平台，更新頻率各異，登山者需人工整合天氣預報、路線狀況、海拔資訊與安全公告，此過程耗時且可能導致不完整或過時的決策。

\subsection*{I. In literature 文獻}
大型語言模型（LLM）在自然語言理解與資訊摘要方面表現強勁，然而在多約束規劃任務上的成功率僅約 0.6\%。這意味著單純依賴生成式 AI 進行登山風險決策並不足夠，LLM 可能產生無來源依據的陳述，或無法識別關鍵的安全約束條件。

\subsection*{J. Judgement 評論}
既有的登山安全資訊系統多為單向資訊發布，缺乏整合性的風險評估與決策支援。現有 LLM 規劃方法在多約束情境下成功率低，且缺乏證據來源追溯與硬性安全紅旗機制，難以直接應用於安全關鍵領域的決策支持。

\subsection*{K. Keypoint 設計精神與 novelty}
TrailGuard TW 的設計精神在於「證據為基礎」與「人類最終決策」。LLM 負責語意理解與解釋，Python 引擎負責定量風險計算，Evidence Ledger 確保每項數據皆可稽查來源，硬性紅旗直接攔截高風險方案，最終決策權保留在人類使用者手中。

\subsection*{L. List the organization 文章結構}
本文依序介紹相關工作（II）、提出之 TrailGuard TW 架構與方法（III）、模擬實驗設計與評估指標（IV）、以及結論與未來工作（V）。

\section{Related Work}
本章節依 M–O 結構，歸納既有方法與本研究的差異。

\subsection*{M. Methods 文獻中有關此問題的解（分類）}
既有高山風險評估方法可分為三類：基於規則的系統、基於資料的統計模型，以及基於 AI 的建議系統。從舊到新、從簡單到複雜進行分類，可發現早期系統多靠靜態規則判斷，較新的系統則嘗試整合即時天氣與路線資料，但多缺乏證據追溯機制。

\subsection*{N. New proposed 新趨勢}
本研究提出結合 LLM 語意理解與 Python 定量引擎的證據導向代理人架構，為既有方法帶來新趨勢。透過時序知識圖譜管理資訊效期，並以 Evidence Ledger 記錄每項建議的來源，系統能在動態變化的登山環境中提供更可靠的決策支援。

\subsection*{O. Organize 歸納}
依 Judgement 評論各類方法的優缺點，以表格歸納既有方法與本研究的差異，如下所示：

\begin{table}[htbp]
\centering
\caption{既有方法與 TrailGuard TW 的比較（4$\times$5 歸納表）}
\label{tab:comparison}
\begin{tabular}{l|p{3cm}|p{3cm}|p{3cm}}
\toprule
評比維度 & 規則式系統 & LLM 僅用系統 & TrailGuard TW \\
\midrule
資訊整合 & 低 & 中 & 高（時序知識圖譜） \\
證據追溯 & 無 & 無 & 有（Evidence Ledger） \\
安全紅旗 & 有（靜態） & 無 & 有（硬性約束） \\
決策權歸屬 & 系統 & 系統 & 人類最終決策 \\
\bottomrule
\end{tabular}
\end{table}

\section{Proposed Scheme (From PA to PM)}
本章節依 PA–PM 結構，詳細說明 TrailGuard TW 的設計與實作。

\subsection*{PA. Aim = Statement 目標}
建立一個可稽核、由人類做最終決策的高風險領域代理人範式，特別適用於登山等安全關鍵戶外活動的行前風險評估。

\subsection*{PB. Based on / Background / Be focus on / Benefit}
高山型國家公園入園人次減半但事故率倍增，行前決策品質是關鍵缺口。本研究聚焦於多約束登山規劃任務，旨在透過證據為基礎的決策支援降低規劃失誤率。

\subsection*{PC. Cure / Cause 解法及有用的原因}
採證據為基礎的決策支援代理人，透過 LLM 語意理解與 Python 定量計算降低規劃失誤。有用的原因在於結合了生成式 AI 的彈性與定量引擎的嚴謹性，同時以 Evidence Ledger 確保數據可稽查。

\subsection*{PD. Design 架構}
TrailGuard TW 的系統架構結合四大組件：LLM 介面負責使用者查詢理解與自然語言解釋；Python 風險引擎執行定量計算，包括 P50／P90 步程估算、天氣窗格判斷與高山症風險評估；時序知識圖譜儲存登山路線、天氣條件、安全公告及其時間屬性；Evidence Ledger 則落實「無來源，無數字」原則，記錄每項建議的證據來源、時間戳與有效期限，硬性紅旗機制直接輸出 NO-GO，防止高風險方案被誤判。

\subsection*{PE. Paper element 元件}
系統主要元件包含：語言模型組件、定量計算引擎、知識圖譜組件、證據帳本組件、安全紅旗組件，以及使用者介面組件。

\subsection*{PF. Formulation / Foundation 功能}
系統提供 P50／P90 步程估算、天氣窗格判斷、高山症風險評估、證據時間戳驗證與硬性安全約束檢查等核心功能。

\subsection*{PG. Graph 架構圖說明}
系統架構圖如下（TikZ 示意）：

\begin{figure}[htbp]
\centering
\begin{tikzpicture}[
    node distance=1.2cm,
    block/.style={
        rectangle,
        draw,
        rounded corners,
        align=center,
        minimum width=2.8cm,
        minimum height=0.9cm,
        font=\footnotesize
    },
    arrow/.style={-{Stealth}, thick}
]
\node[block] (user) {使用者查詢};
\node[block, below=of user] (llm) {LLM 介面};
\node[block, below left=1.2cm and 0.4cm of llm] (engine) {Python 風險引擎};
\node[block, below right=1.2cm and 0.4cm of llm] (graph) {時序知識圖譜};
\node[block, below=1.0cm of llm, yshift=-0.5cm] (ledger) {Evidence Ledger};
\node[block, below=of ledger] (output) {風險評估與建議};

\draw[arrow] (user) -- (llm);
\draw[arrow] (llm) -- (engine);
\draw[arrow] (llm) -- (graph);
\draw[arrow] (engine) -- (ledger);
\draw[arrow] (graph) -- (ledger);
\draw[arrow] (ledger) -- (output);
\end{tikzpicture}
\caption{TrailGuard TW 系統架構圖}
\label{fig:architecture}
\end{figure}

\subsection*{PH. How 達成設計目標}
透過 LLM 解析使用者需求，Python 引擎計算多約束風險，Evidence Ledger 追溯數據來源，硬性紅旗機制攔截高風險方案，最終由人類使用者作出決策，達成證據為基礎的高風險決策支援目標。

\subsection*{PI. Implementation / Interaction among elements / flow chart or Pseudo code}
本研究擬採用 Python 實作風險計算引擎，與 LLM 透過 API 互動，時序知識圖譜以圖資料庫儲存，Evidence Ledger 以結構化日誌記錄。元件間互動流程：使用者輸入登山路線與條件 $\rightarrow$ LLM 解析需求 $\rightarrow$ Python 引擎計算 P50／P90 步程與風險指標 $\rightarrow$ 時序知識圖譜提供即時天氣與安全公告 $\rightarrow$ Evidence Ledger 驗證證據效期 $\rightarrow$ 硬性紅旗檢查關鍵風險 $\rightarrow$ 生成證據支持的建議與說明。

\subsection*{PJ. Jump to an example 舉例}
舉例：某登山者規劃前往雪山主峰的單日登山行程，預報有午後雷陣雨與強風預警。TrailGuard TW 經由時序知識圖譜取得即時天氣資訊，Python 引擎計算預估步程與天氣窗格，Evidence Ledger 驗證資訊來源與效期，硬性紅旗機制檢測到強風預警屬於關鍵風險條件，直接輸出 NO-GO，並由 LLM 生成解釋性說明，建議改期或更換路線。

\subsection*{PK. Key points contribution 重點貢獻}
本研究的重點貢獻包括：證據可稽核（Evidence Ledger）、人類最終決策（Human-in-the-loop）、多約束風險定量化（Python 引擎）、以及資訊時效管理（時序知識圖譜）。

\subsection*{PL. Later how to do simulation}
以 120 個登山情境案例進行模擬實驗，涵蓋不同天候、路線難度、海拔與行程長度，透過比較對照組與消融組驗證系統效能。

\subsection*{PM. Math proof 數學證明}
錯誤安全率（FSR）定義為：

\begin{equation}
FSR = \frac{N_{\text{unsafe} \rightarrow \text{safe}}}{N_{\text{unsafe}}}
\end{equation}

其中 $N_{\text{unsafe} \rightarrow \text{safe}}$ 為不安全情境被錯誤分類為安全的數量，$N_{\text{unsafe}}$ 為實際不安全情境的總數。FSR 越低，代表系統安全性越高。

\section{Simulation}
本章節依 Q–V 結構，說明模擬實驗設計。

\subsection*{Q. Quality and Quantity 實驗量與性質指標}
模擬實驗以 120 個登山情境案例進行，核心性質指標為錯誤安全率（FSR），量化指標包括步程估算誤差（MAPE）、證據引用成功率（CSR）等。

\subsection*{R. Related method to compare 比較方法}
比較對照組為傳統規則式登山建議系統，消融組分別移除 Python 風險引擎、時序有效性檢查、Evidence Ledger 與硬性安全約束四個組件，以驗證各組件的貢獻。

\subsection*{S. Simulation Setup 模擬設置}
模擬情境涵蓋不同天候條件（晴、雨、霧、強風）、路線難度（易、中、難）、海拔高度（2000–3500 公尺）與行程長度（半日、單日、多日），參數調整包括 LLM 提示詞模板、風險閾值與紅旗觸發條件。

\subsection*{T. Tuning 參數調整}
透過網格搜尋與交叉驗證調整風險閾值，使 FSR 最小化的同時保持適當的靈敏度；LLM 提示詞經過多輪迭代優化，確保輸出解釋與定量計算結果一致。

\subsection*{U. Useful results 有用結果}
預期結果顯示 TrailGuard TW 在 FSR、CSR 與 MAPE 三項指標上均優於對照組與消融組，特別在證據引用成功率與風險識別準確度方面表現顯著。

\subsection*{V. Verify 驗證有效性}
透過消融組實驗驗證各組件對系統可靠性的貢獻，證明 TrailGuard TW 能有效降低規劃失誤並提升決策的可稽核性。

\section{Conclusion}
本章節依 W–Z 結構，總結研究意義與未來工作。

\subsection*{W. What proposed 回顧研究意義、動機、挑戰}
本研究回應「開放山林」後高山事故倍增與行前決策品質缺口的課題，提出以證據為基礎的高山風險決策支援代理人 TrailGuard TW。

\subsection*{X. eXcel 專精情境}
系統專精於高風險戶外活動規劃，特別適用於多約束登山決策情境，演算法優點在於結合生成式 AI 與定量風險計算，並以 Evidence Ledger 確保證據可稽查。

\subsection*{Y. Yields 產出}
實驗數據與模擬結果顯示以錯誤安全率（FSR）為核心指標之成效，預期建立可稽核、由人類做最終決策的高風險領域代理人範式。

\subsection*{Z. Zen 結論與未來工作}
最終強調問題獲得解決與研究貢獻：建立一個證據為基礎、人類最終決策的高風險領域代理人。未來工作將著重於完整系統實作、收集真實登山情境資料，並進行更大規模的實驗驗證。

\end{document}
